%% kapa: exam=1 round=27 year=2016 period=1 form=A subject=real_estate subject_ko=부동산학원론 seq=3 number=83
%% source: Q-Net 감정평가사 기출문제 (cst00309) artlSeq=5169402 file=extracted/1008/5169402/제27회 감정평가사 1차 문제/제27회 감정평가사 1교시문제(A형).pdf
%% conversion: kordoc 4.24.51 -> tools/gen1.py (auto); review=proofread-sonnet; flags=none
\begin{question}{83}
다음은 토지에 관하여 설명한 내용들이다. 옳은 것을 모두 고른 것은?
\begin{framedtext}
ㄱ. 택지는 토지에 건물 등의 정착물이 없고 공법이나 사법의 제한을 받는 토지를 말한다. \par
ㄴ. 획지는 법률상의 단위개념으로 소유권이 미치는 범위를 말한다. \par
ㄷ. 이행지는 용도적 지역의 분류 중 세분된 지역 내에서 용도에 따라 전환되는 토지를 말한다. \par
ㄹ. 후보지는 임지지역, 농지지역, 택지지역 상호 간에 다른 지역으로 전환되고 있는 지역의 토지를 말한다. \par
ㅁ. 건부지는 관련법령이 정하는 바에 따라 재난시 피난 등 안전이나 일조 등 양호한 생활환경 확보를 위해, 건축하면서 남겨놓은 일정면적 부분의 토지를 말한다. \par
\end{framedtext}
\begin{choices}
  \item ㄷ
  \item ㄱ, ㄴ
  \item ㄷ, ㄹ
  \item ㄱ, ㄹ, ㅁ
  \item ㄴ, ㄷ, ㄹ
\end{choices}
\end{question}
\answer{3}
